\documentclass[11pt]{article}

\usepackage[T1]{fontenc}
\usepackage[utf8]{inputenc}
\usepackage{lmodern}
\usepackage{amsmath, amssymb}
\usepackage{geometry}
\geometry{margin=1in}

\title{A Simple One-Parameter ``Time-to-Build'' Model of Regulatory Stringency/Uncertainty}
\author{}
\date{}

\begin{document}
\maketitle

\begin{itemize}
  \item \textbf{Demand.} Let the local (rental or price) level be $p(Q)$ with $p'(Q)<0$, where $Q$ is total housing supply delivered.

  \item \textbf{Developer choice (density).} A project on a site chooses units $q\ge 0$ (density) with cost
  \[
    C(q)=F+cq+\tfrac{\gamma}{2}q^2 \quad (\gamma>0).
  \]

  \item \textbf{Regulatory stringency/uncertainty as a single effective delay.} A local regulatory environment $s$ maps into one \textbf{effective time-to-build} (entitlement + construction pipeline) $T(s)$ with $T'(s)>0$. We interpret $T(s)$ broadly: it is a sufficient statistic for any combination of delay, redesign, negotiation, financing frictions, and regulatory uncertainty that pushes cashflows further into the future (and/or reduces their expected present value).

  If revenue arrives only upon completion, expected discounted revenue is summarized by
  \[
    e^{-rT(s)},
  \]
  where $r$ is the relevant discount/financing rate and $T(s)$ is the effective delay. In the data, $T(s)$ can be proxied by observed time-to-build (e.g., permit-to-completion or permit-to-delivery), recognizing it is an imperfect but interpretable summary of the regulatory wedge.

  \item \textbf{Developer problem.} Taking $p$ as given at the project level,
  \[
    \max_{q\ge 0}\;\; e^{-rT(s)}\,p\,q - C(q).
  \]
  The first-order condition (interior) gives
  \[
    q^*(T(s),p)=\frac{e^{-rT(s)}p - c}{\gamma}
    \;\;\text{if }e^{-rT(s)}p>c,\;\text{else }q^*=0.
  \]
  So $\partial q^*/\partial T(s)<0$: longer/more burdensome processes reduce density and can flip projects from ``build'' to ``don't build.''

  \item \textbf{Market equilibrium.} Aggregate supply is $Q=\sum_i q_i^*$. Higher $T(s)$ lowers $Q$; with $p'(Q)<0$, equilibrium prices/rents rise.
\end{itemize}

\section*{Empirical predictions (reduced form)}
\begin{enumerate}
  \item \textbf{Density:} Longer time-to-build $\Rightarrow$ lower completed density $q$ (and more zeros: fewer projects proceed).
  \item \textbf{Quantities:} Longer time-to-build $\Rightarrow$ fewer starts/completions (lower $Q$).
  \item \textbf{Prices/rents:} Longer time-to-build $\Rightarrow$ higher equilibrium $p$.
  \item \textbf{Stronger when financing is tighter:} Effects scale with $r$ (higher rates amplify the penalty of delay).
  \item \textbf{Selection:} Completed projects may be ``higher-return'' on observables (since only those with high $p$ survive), even as overall supply falls.
\end{enumerate}

\section*{Why ``time to build'' can be a sufficient proxy (even if incomplete)}
In this stylized setup, the regulatory environment affects outcomes through the single sufficient index $rT(s)$: regulation acts like a multiplicative wedge $e^{-rT(s)}$ on the present value of revenue. This is exactly the object that matters for the developer's density choice and build/no-build margin. Other regulatory dimensions (fees, design mandates) could be added as additional costs $C(\cdot)$; omitting them does not change the core logic that longer observed time-to-build is a direct, interpretable proxy for the \emph{effective} regulatory wedge that lowers density, reduces aggregate supply, and (in equilibrium) raises prices/rents.

\end{document}
